% ============================================================
%  config.tex — Packages, couleurs, commandes globales
% ============================================================

\usepackage[a3paper, landscape, margin=6mm]{geometry}
\usepackage[T1]{fontenc}
\usepackage[utf8]{inputenc}
\usepackage[french]{babel}
\usepackage{lmodern}
\usepackage{microtype}

% Mise en page
\usepackage{multicol}
\usepackage{array}
\usepackage{tabularx}
\usepackage{booktabs}
\usepackage{colortbl}

% Graphiques / boîtes
\usepackage{tcolorbox}
\tcbuselibrary{skins, most}
\usepackage{tikz}

% Typographie
\usepackage{fontawesome5}
\usepackage{xcolor}
\usepackage{enumitem}
\usepackage{amssymb}

% Entête-pied
\usepackage{fancyhdr}
\usepackage{lastpage}

% ============================================================
%  PALETTE DE COULEURS
% ============================================================
\definecolor{pse-rouge}    {HTML}{C0392B}
\definecolor{pse-bleu}     {HTML}{1A5276}
\definecolor{pse-vert}     {HTML}{1E8449}
\definecolor{pse-orange}   {HTML}{D35400}
\definecolor{pse-jaune}    {HTML}{F39C12}
\definecolor{pse-gris}     {HTML}{717D7E}%{616161}
\definecolor{pse-ligne}    {HTML}{BDC3C7}
\definecolor{pse-titre-bg} {HTML}{2C3E50}

% ============================================================
%  COMMANDES UTILITAIRES
% ============================================================

% Hauteur fixe d'un champ d'acronyme (2 lignes)
\newlength{\fieldheight}
\setlength{\fieldheight}{10.5mm}

% Champ d'acronyme 1 ligne (ref à droite en gris)
\newcommand{\acronymfieldsingle}[4]{%
  \par\noindent
  \textbf{\textcolor{#1}{\large #2}}%
  \rlap{\textcolor{#1!45}{\small #3}}%
  {\color{pse-ligne}\dotfill}%
  {\textcolor{pse-gris!40}{\small #4}}%
  \par\vspace{1mm}%
}

% Champ d'acronyme compact 2 lignes (4 arguments)
%   #1 = couleur de zone
%   #2 = lettre(s) majuscule(s) (grande, grasse, colorée)
%   #3 = suite du mot (collée à la lettre, filigrane sombre)
%   #4 = valeurs de référence / options (ligne 2, gros filigrane ; vide = pointillés seuls)
\newcommand{\acronymfield}[4]{%
  \par\noindent
  \begin{minipage}[c][\fieldheight][t]{\linewidth}%
    \textbf{\textcolor{#1}{\large #2}}%
    \rlap{\textcolor{#1!45}{\small #3}}%
    {\color{pse-ligne}\dotfill}\par\nobreak
    \noindent
    {\color{pse-ligne}\dotfill}%
    {\textcolor{pse-gris!40}{\small #4}}%
  \end{minipage}%
  \par\vspace{1mm}%
}

% Variante : mot en grand sur les 2 lignes de pointillés
%   Le texte (lettre + suite) est doublé en taille et centré verticalement
%   sur toute la hauteur du champ, par-dessus les pointillés.
\newcommand{\acronymfieldbig}[4]{%
  \par\noindent
  \begin{minipage}[c][\fieldheight][t]{\linewidth}%
    \rlap{\smash{\raisebox{-3.5mm}{\textbf{\textcolor{#1}{\Large #2}}%
      \textcolor{#1!40}{\large #3}}}}%
    {\color{pse-ligne}\dotfill}\par\nobreak
    \noindent
    {\color{pse-ligne}\dotfill}%
    {\textcolor{pse-gris!40}{\small #4}}%
  \end{minipage}%
  \par\vspace{1mm}%
}

% Taille intermédiaire : scriptsize + 0.5pt (8.5pt au lieu de 8pt)
\newcommand{\smallscript}{\fontsize{8.5}{10}\selectfont}

% Rond numéroté pour échelle de sévérité (0–10)
\newcommand{\sev}[1]{%
  \tikz[baseline=-0.5ex]\node[draw=pse-jaune!50, circle, minimum size=5mm, inner sep=0pt, font=\smallscript]{#1};\!%
}

% Helper : option à cocher (usage : \opt{Normale} \opt{Superficielle} …)
\newcommand{\opt}[1]{%
  $\square$\,{\small #1}\enspace
}

% Variante avec cases à cocher sur la ligne 2
%   #4 = contenu libre (utiliser \opt{…} pour chaque option)
\newcommand{\acronymfieldcheck}[4]{%
  \par\noindent
  \textbf{\textcolor{#1}{\large #2}}%
  \rlap{\textcolor{#1!45}{\small #3}}%
  {\color{pse-ligne}\dotfill}\;#4%
  \par\vspace{1mm}%
}

% Variante big + cases à cocher
\newcommand{\acronymfieldbigcheck}[4]{%
  \par\noindent
  \rlap{\smash{\raisebox{-1mm}{\textbf{\textcolor{#1}{\Large #2}}%
    \textcolor{#1!40}{\large #3}}}}%
  {\color{pse-ligne}\dotfill}\;#4%
  \par\vspace{1mm}%
}

% Ligne d'écriture simple
\newcommand{\writeline}[1]{%
  \noindent\textbf{\small #1}\enspace{\color{pse-ligne}\dotfill}\par\vspace{2.5pt}%
}

% Sous-titre de bloc
\newcommand{\blocktitle}[2]{%
  \vspace{1.5mm}%
  \noindent{\colorbox{#1}{\color{white}\bfseries\scriptsize\,#2\,}}%
  \vspace{1.5mm}\par
}

% ============================================================
%  EN-TÊTE / PIED DE PAGE
% ============================================================
\pagestyle{fancy}
\fancyhf{}
\renewcommand{\headrulewidth}{0pt}
\renewcommand{\footrulewidth}{0.4pt}
\fancyfoot[L]{\scriptsize\textcolor{pse-gris}{Bilan complémentaire — 4\textsuperscript{e} regard}}
\fancyfoot[C]{\scriptsize\textcolor{pse-gris}{\textit{Fiche plastifiée — feutre effaçable}}}
\fancyfoot[R]{\scriptsize\textcolor{pse-gris}{\thepage/\pageref{LastPage}}}
